\begin{abstract}
Federated learning on wearable and IoT devices spends a scarce resource twice: battery energy when the model is trained, and again every time it is used. We study federated training of multi-exit networks, where a device can train only the first blocks of the network to save energy, and ask how clients, depth, upload precision and local epochs should be chosen under per-device energy budgets. A convergence bound shows that depth heterogeneity creates an error floor that grows with the share of selected clients that skip a block, and gives the coverage needed for a target accuracy. We turn this into EcoFed, a drift-plus-penalty controller with per-device energy queues, a shadow energy price and block-coverage queues, which costs $O(N|\mathcal C|)$ per round. On two public wearable benchmarks (UCI HAR, PAMAP2; subject-wise splits, 16 paired seeds, hyper-parameters tuned on separate seeds) and with an assumed (not measured) energy model, EcoFed raises the final accuracy on strongly label-skewed PAMAP2 at a binding budget by 10 to 27 points over seven baselines (significant after Holm correction), and by 12 points over a static capability-based depth rule. On near-IID UCI HAR it stays within 1 to 3 points of the best baselines while using 1.6 to 2.2 times less energy and leaving 0.8--2.3\% instead of 17--31\% of devices drained. Ablations show that the energy queue is essential and that the coverage queue protects the deep exits (up to 21 points), a controlled probe confirms the predicted rise of the error with missing coverage, and the advantage grows when radio energy is expensive. A lifecycle analysis shows that confidence-gated exits cut inference energy 2 to 7 times and that serving overtakes training energy within two days. EcoFed does not dominate: a static depth rule is better at the tightest budget when training can be stopped at the best checkpoint, and all energy figures are model-based, not measured.
\end{abstract}
\begin{keyword}
Energy-aware AI \sep federated learning \sep multi-exit networks \sep adaptive inference \sep edge computing \sep wearable sensing
\end{keyword}
